\documentclass[twocolumn]{article}

\usepackage[utf8]{inputenc}
\usepackage[T1]{fontenc}
\usepackage{amsmath, amssymb, amsfonts}
\usepackage{graphicx}
%\usepackage{hyperref}
\usepackage[hidelinks]{hyperref}
\usepackage{cite}
\usepackage{booktabs}
\usepackage{float}
\usepackage{caption}
\usepackage{subcaption}

% arXiv metadata
\hypersetup{
    pdftitle={$\Theta$-H Control: From Explosive Divergence to Quasi-Periodic Order in Quartic Nonlinear Systems},
    pdfauthor={Marconi Kolm Madrid},
    pdfsubject={Nonlinear Control, Quasi-Periodic Systems},
    pdfkeywords={$\Theta$-H Control, Quartic Nonlinearity, Quasi-Periodic Systems, Explosive Dynamics, Lyapunov Exponent},
    pdfpagelabels=true,
    plainpages=false,
}

% Page settings
\setlength{\columnsep}{0.5cm}
\setlength{\parindent}{0cm}
\setlength{\parskip}{0.3cm}
\setlength{\textwidth}{6.5in}
\setlength{\textheight}{9in}
\setlength{\oddsidemargin}{0in}
\setlength{\topmargin}{0in}

% Keywords command
\providecommand{\keywords}[1]{\textbf{Keywords:} #1}

% Title
\title{$\Theta$-H Control: From Explosive Divergence to Quasi-Periodic Order\\
\large{in Quartic Nonlinear Systems}}

\author{
\parbox{0.9\textwidth}
\centering
Marconi Kolm Madrid\\
\texttt{madrid@unicamp.br}\\
[0.5em]
Departamento de Sistemas de Energia\\
Faculdade de Engenharia Elétrica e de Computação\\
Universidade Estadual de Campinas (UNICAMP)
}
\date{\today}


\begin{document}

\maketitle

\begin{abstract}
This paper presents $\Theta$-H, a hybrid adaptive controller capable 
of transforming explosive dynamical systems into stable quasi-periodic 
behavior. The application is demonstrated to a system with quartic nonlinearity ($x^4$), which in open loop diverges to infinity in finite time -- a more severe behavior than chaos Under $\Theta$-H 
action, the explosive system is stabilized and organized into a 
quasi-periodic attractor with Lyapunov exponent $\lambda = 0.0014 
\approx 0$, characterized by a torus in the Poincar\'e map and the 
presence of multiple coexisting orbits. The controller operates 
without prior knowledge of the model, learning the dynamics through 
safe autonomous exploration. Experimental results show a spontaneous 
transition between low and high amplitude orbits after approximately 
4300 crossings, maintaining the same fundamental frequency. $\Theta$-H 
represents a new paradigm for controlling systems with severe 
nonlinearities, offering an alternative to traditional mathematical 
modeling.
\end{abstract}

\keywords{$\Theta$-H Control, Quartic Nonlinearity, Quasi-Periodic 
Systems, Explosive Dynamics, Lyapunov Exponent.}

\section{Introduction}

\subsection{Context and Motivation}

Nonlinear dynamical systems exhibit complex behaviors such as 
bifurcations, chaos, and in extreme cases, explosive divergence 
\cite{Ueda1979, Lorenz1963, Ott1990}. While chaos has been extensively 
studied, systems with higher-order nonlinearities (quartic, quintic) 
remain less explored due to their tendency toward catastrophic 
instability \cite{Khalil2002}.

The Duffing oscillator ($x^3$) \cite{Ueda1979} is a paradigmatic 
example of chaos. However, systems with $x^4$ nonlinearity -- here 
termed quartic oscillators -- exhibit qualitatively different behavior: 
explosive divergence in open loop. Controlling such systems is a 
significantly greater challenge, as there is no natural attractor 
to tame.

Recent advances in data-driven control have introduced 
model-free approaches based on machine learning \cite{Koeglmayr2026}, 
but these require extensive offline training. Adaptive methods 
for high-order nonlinear systems \cite{Zhang2025} achieve 
fixed-time convergence at the cost of high computational 
complexity. Chen et al. \cite{Chen2007} demonstrated the 
generation of quasi-periodic and chaotic attractors via 
switching schemes, though without autonomous learning 
capabilities. The present work fills this gap with Theta-H, 
a controller that combines computational simplicity, 
real-time learning, and the unique ability to handle 
explosive systems.

\subsection{Contributions}

This paper presents the following contributions:

\begin{enumerate}
    \item A novel hybrid adaptive controller, called $\Theta$-H,
    combining state feedback with BK memory and autonomous learning.
    
    \item The first demonstration of controlling an explosive 
    quartic system, transforming its dynamics into stable 
    quasi-periodic behavior.
    
    \item Discovery of a rare phenomenon: multiple coexisting 
    quasi-periodic orbits on the same torus, with spontaneous 
    transition between low and high amplitude regimes.
    
    \item Experimental evidence of Lyapunov exponent 
    $\lambda = 0.0014 \approx 0$, confirming the quasi-periodic 
    nature of the controlled system.
\end{enumerate}

\subsection{Paper Organization}

Section~\ref{sec:theta_h} describes the $\Theta$-H architecture.
Section~\ref{sec:mechanics} details the internal mechanics, including 
the BK memory, safe autonomous exploration.
Section~\ref{sec:system} presents the quartic system. 
Section~\ref{sec:results} shows experimental results. 
Section~\ref{sec:discussion} discusses implications. 
Section~\ref{sec:conclusion} concludes.

\section{$\Theta$-H Architecture}
\label{sec:theta_h}

\subsection{General Structure}

$\Theta$-H is defined by:
\begin{equation}
u(t) = -K \cdot x(t) + u_\Theta(t) + u_{BK}(t)
\label{eq:theta_h}
\end{equation}
where:
\begin{itemize}
    \item $K = [K_1, K_2]$ are fixed gains ($K_1=8$, $K_2=6$ in this work)
    \item $u_\Theta$(t) is the adaptive (learned) action
    \item $u_{BK}(t)$ is the BK memory (low-pass filter)
\end{itemize}

\subsection{$\Theta$ Action}

The $\Theta$ action is given by an adaptive PID controller:
\begin{equation}
u_\theta = K_p e + K_i \int e \, dt + K_d \dot{e}
\label{eq:theta_action}
\end{equation}
With gains derived from $\tau$ parameters:
\begin{equation}
K_p = \frac{1.2}{\tau_1}, \quad 
K_i = \frac{0.6}{\tau_1 \tau_2}, \quad 
K_d = 0.5\, \tau_1\, \tau_3
\label{eq:gains}
\end{equation}
In this work, the values $\tau_1 = 0.5$, $\tau_2 = 0.4$, $\tau_3 = 0.1$ are used.

\subsection{BK Memory}

The BK memory acts as a low-pass filter:
\begin{equation}
u_{BK}(t) = \alpha \cdot u_{BK}(t-dt) + (1-\alpha) \cdot u_{fb}(t)~,
\label{eq:bk}
\end{equation}
where $\alpha = 0.98$ is the memory factor, and $u_{fb}(t) = -K \cdot x(t)$.

\subsection{Autonomous Learning}

$\Theta$-H learns the plant dynamics through:

\begin{enumerate}
    \item Safe exploration with increasing amplitude impulses
    \item Identification of the growth rate $\lambda$
    \item Autonomous adjustment of the time window
\end{enumerate}

\section{$\Theta$-H Internal Mechanics: The Learning Engine}
\label{sec:mechanics}

\subsection{The Memory-BK Low-Pass Filter}

One of the key innovations of $\Theta$-H is the BK memory mechanism:
{\small
\begin{equation}
u_{BK}(t) = \alpha u_{BK}(t-dt) + (1-\alpha) u_{fb}(t), \quad \alpha = 0.98
\label{eq:bk_filter}
\end{equation}
}
This first-order low-pass filter introduces a time-delayed feedback 
that creates an emergent quasi-periodic structure. Without BK ($\alpha=0$), 
the controller becomes purely instantaneous and the system remains chaotic 
or explosive. With BK ($\alpha \approx 1$), the controller gains a 
memory of its own past actions, creating a self-referential 
dynamical system.

\subsection{Safe Autonomous Exploration of Explosive Plants}

Learning the dynamics of an explosive system requires extreme caution. 
$\Theta$-H employs a hierarchical exploration strategy:

\begin{enumerate}
    \item Impulse Testing: Small amplitude impulses ($A = 0.01$) 
    are applied. If the response is below a safety threshold ($|y| < 0.1$),
    amplitude is doubled.
    
    \item Time Window Adaptation: The identification window is 
    automatically shifted later in time ($t_{window} = 1.2, 1.4, 1.6, ...$) 
    until the estimated growth rate $\lambda$ stabilizes.
    
    \item Learning Convergence: After approximately 10-15 safe 
    impulses, the system's effective parameters are identified with 
    sufficient accuracy to design the compensator.
\end{enumerate}

\subsection{Why $\Theta$-H Transforms Explosive Dynamics into Quasi-Periodicity}

The controlled quartic system exhibits an emergent property: 
spontaneous transitions between low and high amplitude orbits 
with identical fundamental frequency. This phenomenon is explained by 
the interaction between:

\begin{enumerate}
    \item The linearized plant dynamics (after compensation)
    \item The BK memory (introducing phase delay)
    \item The residual nonlinearity (uncompensated $x^4$ term)
\end{enumerate}

Mathematically, the closed-loop system approximates:
\begin{equation}
x'' + \delta x' + \beta x = F\cos(\omega t) + \varepsilon(x,t)~,
\label{eq:closed_loop}
\end{equation}
where $\varepsilon(x,t)$ is the residual error from imperfect compensation. 
This residual, though small, is sufficient to maintain the system on a 
torus rather than collapsing to a point or diverging.

\subsection{Comparison with Existing Methods}

\begin{table*}[htb]
\centering
\caption{Comparison of Control Methods for Explosive Nonlinear Systems}
\label{tab:comparison_methods}
{\small
\begin{tabular}{|l|c|c|c|c|}
\hline
\textbf{Method} & \textbf{Model Required} & \textbf{Adaptive} & \textbf{Explosive Plants} & \textbf{Quasi-Periodic Output} \\
\hline
PID & No & No & No & No \\
\hline
LQR & Yes & No & No & No \\
\hline
Sliding Mode & Partial & Yes & Limited & No \\
\hline
MPC & Yes & Yes & No & No \\
\hline
Backstepping & Yes & No & No & No \\
\hline
$\Theta$-H & No & Yes & Yes & Yes \\
\hline
\end{tabular}
}
\end{table*}

\subsection{The Emergent Torus: A New Type of Attractor}

Perhaps the most surprising result is that $\Theta$-H does not force the 
system to a fixed point, but rather to a quasi-periodic torus. 
This is counterintuitive: a controller designed to stabilize might be 
expected to drive the system to equilibrium. However, the combination 
of state feedback ($-Kx$) and BK memory creates a self-sustaining 
oscillation at a frequency determined by the closed-loop dynamics.

The torus is characterized by:
\begin{equation}
r(\theta) = r_0 + A_1 \sin(2\pi \theta / T) + A_2 \sin(4\pi \theta / T) + \cdots
\label{eq:torus}
\end{equation}
where $r_0 \approx 0.384$, $T \approx 147$ crossings, and the modulation 
amplitudes $A_1 \approx 0.155$, $A_2 \approx 0.078$.

Table~\ref{tab:torus_params} summarizes the emergent torus characteristics.

\begin{table}[htb]
\centering
\caption{Emergent Torus Characteristics}
\label{tab:torus_params}
{\small
\begin{tabular}{|l|c|}
\hline
\textbf{Parameter} & \textbf{Value} \\
\hline
Mean radius $r_0$ & 0.384 \\
\hline
Radius standard deviation & 0.155 (40\%) \\
\hline
Angular increment & -0.0426 rad \\
\hline
Points per revolution & $\approx 147.5$ \\
\hline
Fundamental period & $\approx 128$ s \\
\hline
Low amplitude regime & $\pm 0.1$ \\
\hline
High amplitude regime & $\pm 0.6$ \\
\hline
Transition point & $\approx 4300$ crossings \\
\hline
Lyapunov exponent & $0.0014 \approx 0$ \\
\hline
Attractor center & $(-0.0102, -0.0102)$ \\
\hline
\end{tabular}
}
\end{table}

\subsection{Energy Considerations}

The observed spontaneous transition from low to high amplitude ($\pm 0.1$ 
to $\pm 0.6$) occurs after $\approx 4300$ crossings ($\approx 3600$ s). 
This suggests the existence of multiple stable orbits on the same 
torus, and the system can "jump" between them when perturbed by numerical 
noise or residual errors. A similar structure of concentric ellipses 
in phase space was observed by Chen et al. \cite{Chen2007} using 
switching schemes, suggesting that toroidal attractors may be a 
universal feature of certain controlled systems.

The energy required for control is reflected in the magnitude of 
$u_{BK}(t)$, which reaches peaks of $\approx 10$ during transitions.

\subsection{What Makes $\Theta$-H Different}

Unlike conventional controllers that attempt to eliminate 
nonlinearities, $\Theta$-H compensates them just enough to make 
the system behave as desired. It does not require:

\begin{itemize}
    \item Mathematical model of the plant
    \item Prior knowledge of the nonlinearity order ($x^3$, $x^4$, or higher)
    \item Gain scheduling or parameter tuning
    \item Separation of identification and control phases
\end{itemize}

The controller learns the necessary compensation through safe 
exploration and stores it implicitly in the BK memory and integral 
action. This is a fundamental departure from model-based control.

\section{Quartic System}
\label{sec:system}

\subsection{Definition}

The following system is considered:
\begin{equation}
x'' + \delta x' + \beta x + \alpha x^4 = F \cos(\omega t) + u
\label{eq:quartic}
\end{equation}
With parameters: $\delta = 0.25$, $\beta = -1$, $\alpha = 1$, $F = 0.3$, 
$\omega = 1.2$.

\subsection{Open Loop Behavior}

Without control ($u=0$), the system diverges to $\pm\infty$ in finite 
time. This contrasts with the Duffing oscillator ($x^3$), which remains 
bounded even in chaotic regimes.

\subsection{Control Challenge}

Controlling this system is particularly challenging because:

\begin{enumerate}
    \item There is no natural attractor to stabilize
    \item The quartic nonlinearity grows rapidly
    \item Small control errors can lead to explosion
\end{enumerate}

\section{Experimental Results}
\label{sec:results}

\subsection{Experimental Setup}

Numerical simulation was performed with:
\begin{itemize}
    \item Integration step: $dt = 0.001$ s
    \item Total time: $20,000$ s
    \item $\Theta$-H controller with $\tau_1=0.5$, $\tau_2=0.4$, $\tau_3=0.1$
    \item Gains $K$: $K_1=8$, $K_2=6$
\end{itemize}

\subsection{Poincar\'e Map}

The Poincar\'e map of the controlled system, shown in Figure~\ref{fig:poincare}, 
is constructed from crossings of the section defined by $x_2 = 0$ with 
increasing phase. This map reveals a striking structure: two concentric 
nearly-circular ellipses centered at $(-0.0102, -0.0102)$, practically 
at the origin.

\begin{figure}[htb]
\centering
\includegraphics[width=0.9\columnwidth]{fig_poincare.pdf}
\caption{Poincar\'e map showing two concentric nearly-circular ellipses. 
The inner ellipse (radius $r_{in} \approx 0.274$) contains $59.6\%$ of the 
points, while the outer ellipse (radius $r_{out} \approx 0.577$) contains 
$40.4\%$. The excentricities are $e_{in} = 0.0392$ and $e_{out} = 0.0057$, 
indicating nearly perfect circles. The radius ratio $r_{out}/r_{in} \approx 
2.11$ reveals a harmonic relationship between the two coexisting orbits.}
\label{fig:poincare}
\end{figure}
The inner ellipse, containing the majority of points ($59.6\%$), corresponds 
to a low-amplitude orbit ($\pm 0.1$), while the outer ellipse ($40.4\%$ of 
points) corresponds to a high-amplitude orbit ($\pm 0.6$). Both ellipses 
exhibit extremely low excentricities ($e_{in} = 0.0392$, $e_{out} = 0.0057$), 
meaning they are nearly perfect circles. This near-circularity suggests 
that the dynamics on each torus are highly regular and close to harmonic 
oscillation.

The radius ratio $r_{out}/r_{in} \approx 2.11$ is particularly noteworthy, 
as it is very close to $2$. This harmonic relationship indicates a possible 
resonance between the two coexisting orbits. The spontaneous transition 
between these orbits, observed after approximately $4300$ crossings, 
suggests that the system can "jump" from one quasi-periodic state to 
another without external perturbation -- a phenomenon rarely documented 
in controlled nonlinear systems.

\subsection{Lyapunov Exponent}

Figure~\ref{fig:lyapunov} presents the Lyapunov exponent calculation.

\begin{figure}[htb]
\centering
\includegraphics[width=0.95\columnwidth]{fig_lyapunov.pdf}
\caption{Lyapunov exponent calculation using the Wolf method \cite{Wolf1985}, implementado conforme a abordagem prática de Rosenstein \cite{rosenstein1993}. 
The experimental data (blue dots) show the evolution of the 
logarithmic distance between neighboring trajectories. 
The linear fit (red line) yields a slope of $\lambda = 0.0014$, 
which is approximately zero. This confirms that the controlled 
system is quasi-periodic, not chaotic (typical chaotic systems 
exhibit $\lambda > 0.1$).}
\label{fig:lyapunov}
\end{figure}

The obtained value is:
\begin{equation}
\lambda = 0.0014 \approx 0
\label{eq:lyapunov_value}
\end{equation}
This value is practically zero, confirming that the controlled 
system is quasi-periodic, not chaotic.

\subsection{Spontaneous Transition}

Figure~\ref{fig:transicao} shows a spontaneous transition between 
two operating regimes.

\begin{figure}[htb]
\centering
\includegraphics[width=1\columnwidth]{fig_transicao.pdf}
\caption{Spontaneous transition between low and high amplitude orbits.}
\label{fig:transicao}
\end{figure}

After approximately 4300 crossings ($\approx 3600$ seconds), the system 
spontaneously transitions from a low amplitude orbit ($\pm 0.1$) to a 
high amplitude orbit ($\pm 0.6$), maintaining the same fundamental 
frequency.

\subsection{Frequency Spectrum}

Figure~\ref{fig:spectrum} shows the power spectrum of the controlled 
system's position signal $x(t)$, calculated using the Welch method 
with a Hanning window.

\begin{figure}[htb]
\centering
\includegraphics[width=1\columnwidth]{fig_spectrum}
\caption{Power spectrum of the controlled system (Hanning window). 
The spectrum exhibits a dominant peak at approximately $0.16$ Hz and 
a secondary peak at $0.80$ Hz. A re-excitation is observed between 
$2.6$ and $5.0$ Hz. The ratio $f_2/f_1 \approx 5$ is harmonically 
related to the radius ratio $r_{out}/r_{in} \approx 2.11$ observed 
in the Poincar\'e map (Figure~\ref{fig:poincare}).}
\label{fig:spectrum}
\end{figure}

\section{Discussion}
\label{sec:discussion}

The results presented in this work demonstrate that $\Theta$-H is capable of 
stabilizing an intrinsically explosive quartic oscillator, organizing it 
into a quasi-periodic attractor with Lyapunov exponent $\lambda = 0.0014 
\approx 0$.

The most surprising finding is the emergence of two 
concentric nearly-circular ellipses in the Poincar\'e map, with 
excentricities $e_{in} = 0.0392$ and $e_{out} = 0.0057$, and a radius 
ratio of approximately $2.11:1$. This indicates the existence of 
two distinct coexisting quasi-periodic orbits on the same torus.

The spontaneous transition between low ($\pm 0.1$) and high ($\pm 0.6$) 
amplitude regimes after approximately $4300$ crossings suggests that 
the system possesses at least two stable orbits on the same invariant 
manifold. The fact that both regimes share the same fundamental 
frequency ($\approx 128$ s) indicates that the transition is not a 
traditional bifurcation, but rather a "jump" between two attractors 
within the same torus.

When compared to existing control methods, $\Theta$-H stands out for its 
ability to handle explosive nonlinearities without requiring a 
mathematical model. Traditional controllers such as PID or LQR either 
fail or require extensive tuning; advanced methods like sliding mode 
or backstepping \cite{Slotine1991} demand partial model knowledge and are not designed 
for intrinsically unstable plants. Recent machine learning approaches 
\cite{Koeglmayr2026} offer model-free control but require extensive 
offline training, while adaptive fixed-time methods \cite{Zhang2025} 
achieve convergence guarantees at the cost of high complexity. 
Table~\ref{tab:comparison_methods} summarizes these differences.

The main limitation of the proposed approach is the need for sufficient 
actuator energy to compensate for the explosive nonlinearity. However, the discovery of 
multiple coexisting quasi-periodic orbits opens new questions about 
the structure of controlled nonlinear systems.

In the present simulations, the control signal $u_{BK}(t)$ reached peaks of 
approximately 10 during transitions, which may be demanding for certain 
hardware implementations. However, for systems where energy is not a 
limiting factor, $\Theta$-H offers a powerful and practical solution.

Another important observation is that the residual nonlinearity 
$\varepsilon(x,t)$ -- arising from imperfect compensation -- plays a 
crucial role in maintaining the system on the torus. If compensation 
were perfect, the system would likely collapse to a fixed point or a 
simple limit cycle. The slight imperfection is what gives rise to the 
rich quasi-periodic structure observed. This suggests that $\Theta$-H does 
not simply "cancel" nonlinearities but rather "sculpts" them into a 
desired form.

Finally, the fact that $\Theta$-H operates without knowledge of whether 
the nonlinearity is $x^3$, $x^4$, or any other function suggests a 
degree of universality. The controller learns the necessary compensation 
from observations alone, making it applicable to a wide range of systems 
with severe nonlinearities.

\subsection{Comparison with Literature}

Table~\ref{tab:comparison} compares the controlled quartic system 
with the classical Duffing oscillator. While the Duffing oscillator 
($x^3$) exhibits bounded chaotic behavior, the quartic system ($x^4$) 
diverges explosively in open loop -- a significantly more severe 
instability. The Lyapunov exponent of the controlled quartic system 
($\lambda = 0.0014 \approx 0$) is orders of magnitude smaller than 
typical values for chaotic Duffing ($\lambda \approx 0.1-0.3$), 
confirming the quasi-periodic nature of the controlled dynamics. 
To the best of our knowledge, this is the first demonstration of 
controlling an explosive quartic system and organizing it into a 
stable quasi-periodic attractor with two concentric nearly-circular 
ellipses.
\begin{table*}[htb]
\centering
\caption{Comparison between Duffing ($x^3$) and Quartic ($x^4$) Systems}
\label{tab:comparison}
{\small
\begin{tabular}{|l|c|c|}
\hline
\textbf{Aspect} & \textbf{Duffing ($x^3$)} & \textbf{Quartic ($x^4$)} \\
\hline
Natural behavior & Chaotic & Explosive \\
\hline
Control challenge & Medium & Very High \\
\hline
Lyapunov exponent & $\approx 0.1-0.3$ & $\approx 0.0014$ \\
\hline
This work & - & $\checkmark$ First demonstration \\
\hline
\end{tabular}
}
\end{table*}
The table highlights the fundamental differences between the two 
systems and the unprecedented nature of our results. The quartic 
system, left uncontrolled, diverges to infinity in finite time -- 
a behavior that cannot be characterized as chaotic in the traditional 
sense. The fact that $\Theta$-H not only stabilizes this explosive 
system but also organizes it into a structured quasi-periodic 
attractor with two concentric nearly-circular ellipses represents 
a significant advance in nonlinear control theory.

\subsection{Universality of the Method}

$\Theta$-H requires no prior knowledge of the model, operating through 
autonomous learning. The controller learns the necessary compensation 
through safe exploration, without needing to know the order of the 
nonlinearity ($x^3$, $x^4$, or higher). This suggests the approach 
can be extended to other classes of systems with severe nonlinearities, 
including higher-order polynomial nonlinearities ($x^5$, $x^6$), 
systems with hysteresis, and even chaotic systems such as Lorenz 
or R\"ossler.

The key features that contribute to this universality are:

\begin{itemize}
    \item Model-free operation: No mathematical model of 
    the plant is required.
    
    \item Autonomous learning: The controller learns the 
    dynamics through safe exploration, without offline training data.
    
    \item BK memory: The low-pass filter creates emergent 
    quasi-periodic structures without requiring precise compensation.
    
    \item Robustness: The controller maintains stability 
    across a wide range of parameters, as demonstrated by the 
    spontaneous transitions between orbits.
\end{itemize}

These characteristics make $\Theta$-H a promising candidate for 
real-world applications where accurate models are unavailable 
or where systems exhibit complex, poorly understood dynamics.

\section{Conclusion and Future Work}
\label{sec:conclusion}
This paper has introduced $\Theta$-H, a hybrid adaptive controller that 
combines state feedback with BK memory and autonomous learning. The 
controller was successfully applied to a quartic oscillator ($x^4$) - 
a system that diverges explosively in open loop.

The most significant findings of this work are twofold. First, $\Theta$-H 
transforms the explosive quartic oscillator into a stable quasi-periodic 
attractor with Lyapunov exponent $\lambda = 0.0014 \approx 0$, confirming 
its non-chaotic nature. Second, and more remarkably, the Poincar\'e map 
reveals two concentric nearly-circular ellipses with excentricities 
$e_{in} = 0.0392$ and $e_{out} = 0.0057$, and a harmonic radius ratio 
of approximately $2.11:1$. This indicates the existence of two distinct 
coexisting quasi-periodic orbits on the same torus -- a structure that, 
to the best of our knowledge, has not been previously documented in 
controlled explosive systems.

The spontaneous transition between low ($\pm 0.1$) and high ($\pm 0.6$) 
amplitude regimes after approximately $4300$ crossings further 
demonstrates the richness of the controlled dynamics. Both regimes 
share the same fundamental frequency ($\approx 128$ s), suggesting 
that the transition is not a traditional bifurcation but rather a 
jump between two attractors within the same invariant manifold.

Several directions for future work remain open. The demonstration 
of arbitrary trajectory tracking -- such as ramps or parabolas -- would 
further establish the controller's capabilities. Extension to other 
explosive systems ($x^5$, $x^6$) and to other chaotic systems 
(Lorenz, R\"ossler, Chua) would test the universality of the approach. 
Hardware implementation on a physical circuit would validate practical 
feasibility. A theoretical analysis of the learning convergence and 
stability guarantees would provide mathematical foundations. Finally, 
a deeper investigation into the spontaneous orbit transitions and the 
emergence of multiple nearly-circular ellipses could reveal new 
insights into the structure of quasi-periodic attractors.

In summary, $\Theta$-H represents a paradigm shift in nonlinear system 
control. Instead of asking ``what is this system?'', it asks 
``what do you want it to be?''. With sufficient actuator energy, 
Theta-H delivers.

Beyond its technical contributions, this work raises fundamental 
questions about the nature of controlled nonlinear systems. The emergence
 of two concentric nearly-circular ellipses with harmonic radius ratio 
 \mbox{$2.11:1$} suggests that $\Theta$-H does not merely suppress 
 nonlinearities but actively organizes them into structured, predictable 
 patterns. This is reminiscent of self-organization phenomena observed in nature, 
 from the formation of spiral galaxies to the synchronization of biological 
 oscillators. The fact that such order can emerge from explosive dynamics 
 through a simple, model-free controller hints at a deeper principle: that 
 complex systems, when appropriately compensated, naturally tend toward 
 harmonic structures. Whether this principle extends to other classes 
 of nonlinearities -- such as quintic or higher-order terms -- remains 
 an open question that merits further investigation.

\subsection{Future Work}

Next steps include:
\begin{enumerate}
    \item Demonstration of arbitrary trajectory tracking
    (ramps, parabolas) in the quartic system
    
    \item Extension to other explosive systems ($x^5$, $x^6$)
    
    \item Hardware implementation (Chua circuit, electromechanical 
    systems)
    
    \item Theoretical analysis of learning convergence
    
    \item Investigation of spontaneous orbit transitions
\end{enumerate}

\section*{Acknowledgments}

This study was financed in part by the Coordenação de Aperfeiçoamento de Pessoal de Nível Superior - Brasil (CAPES) - Finance Code 001.

\bibliographystyle{ieeetr}
\bibliography{references.bib}

\end{document} 